% Options for packages loaded elsewhere
\PassOptionsToPackage{unicode}{hyperref}
\PassOptionsToPackage{hyphens}{url}
%
\documentclass[
]{article}
\usepackage{amsmath,amssymb}
\usepackage{iftex}
\ifPDFTeX
  \usepackage[T1]{fontenc}
  \usepackage[utf8]{inputenc}
  \usepackage{textcomp} % provide euro and other symbols
\else % if luatex or xetex
  \usepackage{unicode-math} % this also loads fontspec
  \defaultfontfeatures{Scale=MatchLowercase}
  \defaultfontfeatures[\rmfamily]{Ligatures=TeX,Scale=1}
\fi
\usepackage{lmodern}
\ifPDFTeX\else
  % xetex/luatex font selection
\fi
% Use upquote if available, for straight quotes in verbatim environments
\IfFileExists{upquote.sty}{\usepackage{upquote}}{}
\IfFileExists{microtype.sty}{% use microtype if available
  \usepackage[]{microtype}
  \UseMicrotypeSet[protrusion]{basicmath} % disable protrusion for tt fonts
}{}
\makeatletter
\@ifundefined{KOMAClassName}{% if non-KOMA class
  \IfFileExists{parskip.sty}{%
    \usepackage{parskip}
  }{% else
    \setlength{\parindent}{0pt}
    \setlength{\parskip}{6pt plus 2pt minus 1pt}}
}{% if KOMA class
  \KOMAoptions{parskip=half}}
\makeatother
\usepackage{xcolor}
\usepackage{longtable,booktabs,array}
\usepackage{calc} % for calculating minipage widths
% Correct order of tables after \paragraph or \subparagraph
\usepackage{etoolbox}
\makeatletter
\patchcmd\longtable{\par}{\if@noskipsec\mbox{}\fi\par}{}{}
\makeatother
% Allow footnotes in longtable head/foot
\IfFileExists{footnotehyper.sty}{\usepackage{footnotehyper}}{\usepackage{footnote}}
\makesavenoteenv{longtable}
\setlength{\emergencystretch}{3em} % prevent overfull lines
\providecommand{\tightlist}{%
  \setlength{\itemsep}{0pt}\setlength{\parskip}{0pt}}
\setcounter{secnumdepth}{-\maxdimen} % remove section numbering
\ifLuaTeX
  \usepackage{selnolig}  % disable illegal ligatures
\fi
\IfFileExists{bookmark.sty}{\usepackage{bookmark}}{\usepackage{hyperref}}
\IfFileExists{xurl.sty}{\usepackage{xurl}}{} % add URL line breaks if available
\urlstyle{same}
\hypersetup{
  hidelinks,
  pdfcreator={LaTeX via pandoc}}

\author{}
\date{}

\begin{document}

\hypertarget{knowledge-distillation-and-deployment-of-a-lightweight-command-parsing-lm-on-a-local-cloud}{%
\section{Knowledge Distillation and Deployment of a Lightweight
Command-Parsing LM on a Local
Cloud}\label{knowledge-distillation-and-deployment-of-a-lightweight-command-parsing-lm-on-a-local-cloud}}

\textbf{Name:} Aryan Kumar\\
\textbf{Roll number:} 230101050\\
\textbf{Course:} CS3081 Cloud Computing\\
\textbf{Repository:}
\url{https://github.com/aryanB1706/cloud_computing_assignment}

\hypertarget{architecture-and-design-decisions}{%
\subsection{1. Architecture and design
decisions}\label{architecture-and-design-decisions}}

The project is a local cloud with four layers. A k3d cluster is created
from \texttt{cluster/k3d.yaml} (one server, one agent, Traefik and
metrics-server kept enabled). SeaweedFS runs in namespace \texttt{data}
and exposes an S3-compatible API. A parallel Kubernetes Indexed Job in
namespace \texttt{pipeline} reads
\texttt{datasets/commands\_pool.jsonl}, labels each command with a
host-run Ollama \texttt{phi3.5}, and writes one raw shard object per
worker. A cleaning step validates the replies and writes
\texttt{datasets/v1/train.jsonl} with a SHA-256 checksum. Training
distils the teacher replies into \texttt{EleutherAI/pythia-70m} and
stores \texttt{model.tar.gz} and \texttt{model.sha256} in the
\texttt{models} bucket. The FastAPI service runs in namespace
\texttt{app}, downloads the model at start-up, and is exposed by a
Traefik Ingress.

The main design choice is to keep weights out of the image and out of
Git. The service image contains only Python and pinned libraries;
\texttt{/readyz} stays 503 until the model has been downloaded,
checksum-verified, extracted, and loaded. This makes cold start
observable and makes a model update a change of
\texttt{MODEL\_URI}/\texttt{MODEL\_VERSION}, not a rebuild. Ollama runs
on the host rather than in the cluster because the teacher is only
needed during Task 1 and host networking avoids a multi-gigabyte GPU
image. The workers reach it through \texttt{host.k3d.internal:11434}.
Every component is stateless except SeaweedFS, so Pod deletion and
rolling updates are safe.

Version pins are recorded in the training configuration and the service,
worker, and training dependency files. Cluster configuration is in
\texttt{k8s/}; host-generated S3 credentials are kept in the git-ignored
\texttt{.env.local}.

\hypertarget{task-1-object-storage-and-parallel-teacher-labelling}{%
\subsection{2. Task 1: object storage and parallel teacher
labelling}\label{task-1-object-storage-and-parallel-teacher-labelling}}

\hypertarget{storage-layout}{%
\subsubsection{2.1 Storage layout}\label{storage-layout}}

SeaweedFS is deployed by \texttt{k8s/seaweedfs.yaml} into namespace
\texttt{data}. Two buckets are created: \texttt{datasets} and
\texttt{models}. \texttt{scripts/init\_storage.py} uploads the provided
\texttt{commands\_pool.jsonl} to \texttt{datasets/commands\_pool.jsonl}.
\texttt{prompts/teacher\_prompt.txt} asks for exactly the four keys and
the eight allowed intents, gives three examples, and is fixed at
\texttt{temperature=0}, \texttt{seed=42}, \texttt{top\_k=1}.

\hypertarget{sharding-and-resumability}{%
\subsubsection{2.2 Sharding and
resumability}\label{sharding-and-resumability}}

\texttt{worker/worker.py} assigns each input row to one shard by its
line-number remainder and writes a JSONL shard containing the id,
command, and raw teacher reply. The Indexed Kubernetes Job uses four
completions and parallelism two. Each of the 3,000 input rows belongs to
exactly one shard, so there are no gaps or repeats.

Resume is built in: before labelling, the worker reads its own existing
shard object into a set of completed ids and skips them. Progress is
written after every reply by default (atomic temp-file + rename), so an
interrupted multi-hour run loses at most one command. The cleaner also
refuses to proceed if ids are repeated across shards;
\texttt{-\/-allow-incomplete} is only for smoke runs.

\hypertarget{cleaning-rules}{%
\subsubsection{2.3 Cleaning rules}\label{cleaning-rules}}

\texttt{worker/clean.py} applies rejection precedence (a)
-\textgreater{} (b) -\textgreater{} (c) -\textgreater{} (d) and counts
each record exactly once:

\begin{itemize}
\item
  \begin{enumerate}
  \def\labelenumi{(\alph{enumi})}
  \tightlist
  \item
    \texttt{teacher\_raw} is not valid JSON or is not an object;
  \end{enumerate}
\item
  \begin{enumerate}
  \def\labelenumi{(\alph{enumi})}
  \setcounter{enumi}{1}
  \tightlist
  \item
    keys are not exactly \texttt{intent}, \texttt{date}, \texttt{time},
    \texttt{place};
  \end{enumerate}
\item
  \begin{enumerate}
  \def\labelenumi{(\alph{enumi})}
  \setcounter{enumi}{2}
  \tightlist
  \item
    \texttt{intent} is not one of the eight allowed values;
  \end{enumerate}
\item
  \begin{enumerate}
  \def\labelenumi{(\alph{enumi})}
  \setcounter{enumi}{3}
  \tightlist
  \item
    a slot is neither null nor a string that appears in the command
    (case-insensitive, whitespace-normalised);
  \end{enumerate}
\item
  duplicates are removed by normalised command text, keeping the
  smallest id.
\end{itemize}

Survivors are sorted by numeric id and written to
\texttt{datasets/v1/train.jsonl}; the SHA-256 line is written to
\texttt{datasets/v1/train.sha256}. The accounting invariant below is
asserted in code:

\begin{verbatim}
input = invalid_json + wrong_keys + bad_intent + bad_slot + duplicates + kept
\end{verbatim}

\begin{longtable}[]{@{}lr@{}}
\toprule\noalign{}
Category & Full run \\
\midrule\noalign{}
\endhead
\bottomrule\noalign{}
\endlastfoot
Input records & 3,000 \\
Invalid JSON & 49 \\
Wrong keys & 4 \\
Bad intent & 0 \\
Bad slot & 118 \\
Duplicates & 47 \\
Kept & 2,782 \\
\end{longtable}

The invariant holds exactly:
\texttt{3000\ =\ 49\ +\ 4\ +\ 0\ +\ 118\ +\ 47\ +\ 2782}. Coverage
verification found 3,000 expected and present IDs, with no missing,
repeated, or unexpected IDs. The cleaned dataset SHA-256 is
\texttt{7e286136aaaaf4fd15e86ad4fe3c449b304182a69f7dd2664a0d729f13d33734}.

\hypertarget{sharding-rationale}{%
\subsubsection{2.4 Sharding rationale}\label{sharding-rationale}}

Four shards are the smallest count that gives two waves at
\texttt{parallelism:\ 2}, so the cluster exercises both parallel workers
and resume without creating idle pods. Sharding by line number, rather
than by id, is exactly what the assignment specifies and is independent
of the order of ids. Because every worker owns a disjoint residue class,
the union is complete without a merge step. The cost is that a worker's
commands are interleaved through the pool, but the worker sorts its own
output by id, so each object is deterministic. Four workers also mean
four objects, which is small enough to inspect by hand and large enough
to show that the cleaning step really reads multiple shards. The local
smoke run used the same code with four shards and one worker at a time,
for 50 commands per shard (200 total), because a full CPU teacher run of
3,000 commands takes several hours; the Kubernetes Job is the intended
full path and is resumable.

\hypertarget{task-2-training-the-student}{%
\subsection{3. Task 2: training the
student}\label{task-2-training-the-student}}

\hypertarget{configuration}{%
\subsubsection{3.1 Configuration}\label{configuration}}

\texttt{training/config.yaml} records the six required items:

\begin{longtable}[]{@{}
  >{\raggedright\arraybackslash}p{(\columnwidth - 2\tabcolsep) * \real{0.5000}}
  >{\raggedright\arraybackslash}p{(\columnwidth - 2\tabcolsep) * \real{0.5000}}@{}}
\toprule\noalign{}
\begin{minipage}[b]{\linewidth}\raggedright
Item
\end{minipage} & \begin{minipage}[b]{\linewidth}\raggedright
Value
\end{minipage} \\
\midrule\noalign{}
\endhead
\bottomrule\noalign{}
\endlastfoot
Random seed & 42 (\texttt{torch.manual\_seed}, NumPy, Python
\texttt{random}; fixed in code) \\
Learning rate & 5e-5 \\
Batch size & 4 (gradient accumulation 4, effective 16) \\
Epochs & 2 \\
Maximum sequence length & 256 \\
Library versions & Python 3.12, torch 2.4.1, transformers 4.44.2,
tokenizers 0.19.1, safetensors 0.4.5, huggingface-hub 0.24.6, numpy
1.26.4, PyYAML 6.0.2, boto3 1.35.24 \\
\end{longtable}

\texttt{training/config\_v2.yaml} changes only the seed (43), epochs
(3), version and loss file. \texttt{make\ train} reads only
\texttt{training/config.yaml}; \texttt{make\ train\ SUBSET=50} reads the
same file and caps the rows. \texttt{training/loss.csv} is written and
flushed after every epoch.

\hypertarget{example-format-and-loss}{%
\subsubsection{3.2 Example format and
loss}\label{example-format-and-loss}}

Each record becomes
\texttt{Command:\ \textless{}text\textgreater{}\textbackslash{}nJSON:\ \textless{}teacher\ reply\textgreater{}}.
The prompt tokens are masked with \texttt{-100}, so the cross-entropy is
computed only on the teacher reply tokens plus EOS; this matches ``the
model learns to produce the teacher's reply, given the command'' and
prevents the 70M model from spending capacity on reproducing the prompt.
The model is \texttt{EleutherAI/pythia-70m}, loaded in float32, trained
with AdamW and a linear warm-up schedule on CPU. The target is
sequence-level: the exact teacher string, not logits.

\begin{longtable}[]{@{}lrlrr@{}}
\toprule\noalign{}
Version & Records & Epoch losses & Seed & Wall time \\
\midrule\noalign{}
\endhead
\bottomrule\noalign{}
\endlastfoot
v1 & 1,500 & 0.360529, 0.087982 & 42 & 403.2 s \\
v2 & 1,500 & 0.390328, 0.102742, 0.047950 & 43 & 787.8 s \\
\end{longtable}

\hypertarget{evaluation}{%
\subsubsection{3.3 Evaluation}\label{evaluation}}

Predictions are generated for the teacher (\texttt{phi3.5}, same prompt,
temperature 0, seed 42), untrained Pythia-70M, and the two students. Raw
outputs and score reports are in \texttt{results/}. The official
\texttt{tools/score.py} computes JSON validity, intent accuracy, slot
accuracies and full-frame accuracy.

\begin{longtable}[]{@{}lrrrrrr@{}}
\toprule\noalign{}
Model & JSON validity & Intent & Date & Time & Place & Full frame \\
\midrule\noalign{}
\endhead
\bottomrule\noalign{}
\endlastfoot
Teacher & 99\% & 74\% & 74\% & 87\% & 79\% & 40\% \\
Untrained Pythia & 0\% & 0\% & 0\% & 0\% & 0\% & 0\% \\
Student v1 & 99\% & 68\% & 75\% & 89\% & 75\% & 35\% \\
Student v2 & 100\% & 71\% & 76\% & 89\% & 77\% & 38\% \\
\end{longtable}

Version 2 improves JSON validity to 100\%, intent accuracy to 71\%, and
full-frame accuracy to 38\%. Its full-frame score is 95\% of the
teacher's 40\%, while requiring one additional CPU training epoch.

\hypertarget{sequence-level-versus-logit-based-distillation}{%
\subsubsection{3.4 Sequence-level versus logit-based
distillation}\label{sequence-level-versus-logit-based-distillation}}

Sequence-level distillation trains the student on the teacher's
generated sequence, treating it as a hard target. The loss is the
ordinary next-token cross-entropy on that sequence. Logit-based
distillation trains the student to match the teacher's full next-token
probability distribution at each position, typically with a
temperature-softened KL divergence. Logit-based distillation is not
possible here because the teacher is served by Ollama, which exposes
only the generated text through \texttt{/api/generate}; it does not
return per-token logits. Even if it did, the teacher and student use
different tokenizers, so their vocabularies would need an alignment
step. A limitation of the sequence-level approach is that it discards
the teacher's uncertainty and cannot teach the student about near-miss
alternatives; with only about 1,500 examples, the student can also
memorise surface forms instead of generalising to unseen commands.

\hypertarget{task-3-inference-service-and-container}{%
\subsection{4. Task 3: inference service and
container}\label{task-3-inference-service-and-container}}

The service is a FastAPI application on port 8000 with
\texttt{POST\ /parse}, \texttt{GET\ /healthz} and \texttt{GET\ /readyz}.
\texttt{/parse} returns only the four-key object and sets
\texttt{X-Model-Version} and \texttt{X-Served-By} (the Pod name from a
downward-API environment variable, with \texttt{HOSTNAME} and
\texttt{socket.gethostname()} as fallbacks). The output validator checks
exact keys, the eight intents, and string-or-null slots. If the model
output is not valid, the service returns HTTP 500 with exactly
\texttt{\{"error":"invalid\_model\_output"\}}; it never returns
malformed JSON with 200. \texttt{FAULT\_INJECT=1} replaces the generated
text with \texttt{not\ json} before validation, which exercises that
path.

At start-up the service downloads
\texttt{s3://models/\textless{}version\textgreater{}/model.tar.gz} using
boto3 against \texttt{S3\_ENDPOINT}, verifies \texttt{model.sha256} if
present, safely extracts the archive, and loads the model and tokenizer.
\texttt{MODEL\_URI}, \texttt{MODEL\_VERSION}, \texttt{S3\_ENDPOINT} and
\texttt{MAX\_NEW\_TOKENS} come from the environment; in Kubernetes the
two S3 credentials come only from the \texttt{model-store-credentials}
Secret. The image is built by \texttt{service/Dockerfile} from
\texttt{python:3.12-slim}, installs the CPU-only PyTorch wheel and
pinned dependencies, and copies no weights. It runs as UID 10001 with a
read-only root filesystem and a writable \texttt{/tmp} emptyDir for the
model cache.

The Docker image built successfully at \textbf{1.63 GB} and ran in
Kubernetes with the model fetched from SeaweedFS, not baked into the
image. A local-process 50-request benchmark measured p50 \textbf{267.48
ms}, p95 \textbf{406.08 ms}, throughput \textbf{3.386 requests/s}, peak
RSS \textbf{546.8 MiB}, and \textbf{50/50 successful requests} (0\%
errors; 14.767 s). The required separate Docker run constrained to 1 CPU
and 1 GiB was not measured; local RSS is not a substitute for a
container cgroup measurement. Generation is serialized by a lock.

\hypertarget{task-4-kubernetes-deployment-scaling-and-reliability}{%
\subsection{5. Task 4: Kubernetes deployment, scaling and
reliability}\label{task-4-kubernetes-deployment-scaling-and-reliability}}

\texttt{k8s/app.yaml} creates a two-replica Deployment in \texttt{app}
with CPU request 250m/limit 1, memory request 512Mi/limit 1Gi, liveness
\texttt{/healthz} and readiness \texttt{/readyz},
\texttt{maxUnavailable:\ 0}, \texttt{maxSurge:\ 1}, a non-root read-only
filesystem, and the Secret-backed S3 credentials.
\texttt{k8s/ingress.yaml} routes \texttt{/parse}, \texttt{/healthz} and
\texttt{/readyz} through Traefik; \texttt{k8s/hpa.yaml} scales 1-2
replicas at 70\% CPU with a 60-second scale-down window.

The manifests parse successfully. Live validation was executed on
2026-11-08 with Docker 29.1.3, a two-node k3d cluster, and SeaweedFS.
Both Kubernetes nodes became Ready, the v1 model archive was uploaded to
the live \texttt{models} bucket, and two application replicas became
Ready on separate nodes. Through a Traefik port-forward,
\texttt{/healthz}, \texttt{/readyz}, and \texttt{/parse} each returned
HTTP 200; ten additional inference requests also returned HTTP 200. The
metrics API was available. A pod sample during the load period reported
992m CPU and 375 MiB memory; an earlier idle observation was about
395--398 MiB per pod. HPA scale-up and scale-down events were observed
against the 70\% CPU target.

The host-side k3d mappings on ports 8080 and 8333 returned empty
responses, so the live checks used supported
\texttt{kubectl\ port-forward} paths. The following additional live
measurements were captured: the v1-only rolling update and rollback
completed, pod recovery took 12 seconds, three cold starts took 11, 11,
and 10 seconds, and a short Ingress load test completed 76 requests with
zero failures (p50 6.0 seconds, p95 6.6 seconds, 1.297 requests/second).
HPA scale-down to one replica and subsequent scale-up were also
observed. The short CPU-bound test sent all successful requests to one
replica, so balanced per-replica request share remains unmeasured. The
v2 artifact upload remains blocked by SeaweedFS reporting
\texttt{no\ free\ volumes\ left}; the v1-only rollout validates
Kubernetes behavior but not a real v1-to-v2 model change.

Local limitation: the cluster is two k3d nodes on one laptop and shares
one CPU and one disk. Scaling to two replicas demonstrates the mechanism
but does not provide machine-level fault tolerance or a realistic
multi-node network. A real limitation of the HPA is that CPU utilisation
includes model download and start-up work; a pod can be scaled before it
is Ready, and the readiness probe is what protects the Service from
traffic. During a rolling update the old pod remains Ready until the new
pod downloads and loads the model, so the update can take as long as a
cold start.

\hypertarget{task-5-subnetting-and-nacl}{%
\subsection{6. Task 5: subnetting and
NACL}\label{task-5-subnetting-and-nacl}}

Requirements are allocated largest first, each immediately after the
previous block, with the smallest prefix that fits:

\begin{longtable}[]{@{}
  >{\raggedright\arraybackslash}p{(\columnwidth - 10\tabcolsep) * \real{0.1667}}
  >{\raggedright\arraybackslash}p{(\columnwidth - 10\tabcolsep) * \real{0.1667}}
  >{\raggedright\arraybackslash}p{(\columnwidth - 10\tabcolsep) * \real{0.1667}}
  >{\raggedright\arraybackslash}p{(\columnwidth - 10\tabcolsep) * \real{0.1667}}
  >{\raggedright\arraybackslash}p{(\columnwidth - 10\tabcolsep) * \real{0.1667}}
  >{\raggedright\arraybackslash}p{(\columnwidth - 10\tabcolsep) * \real{0.1667}}@{}}
\toprule\noalign{}
\begin{minipage}[b]{\linewidth}\raggedright
Subnet
\end{minipage} & \begin{minipage}[b]{\linewidth}\raggedright
Network
\end{minipage} & \begin{minipage}[b]{\linewidth}\raggedright
Prefix
\end{minipage} & \begin{minipage}[b]{\linewidth}\raggedright
Usable host range
\end{minipage} & \begin{minipage}[b]{\linewidth}\raggedright
Broadcast
\end{minipage} & \begin{minipage}[b]{\linewidth}\raggedright
Usable
\end{minipage} \\
\midrule\noalign{}
\endhead
\bottomrule\noalign{}
\endlastfoot
app & 10.20.0.0 & /22 & 10.20.0.1-10.20.3.254 & 10.20.3.255 & 1022 \\
public & 10.20.4.0 & /24 & 10.20.4.1-10.20.4.254 & 10.20.4.255 & 254 \\
data & 10.20.5.0 & /25 & 10.20.5.1-10.20.5.126 & 10.20.5.127 & 126 \\
management & 10.20.5.128 & /27 & 10.20.5.129-10.20.5.158 & 10.20.5.159 &
30 \\
\end{longtable}

The app subnet needs 1,000 hosts, so /22 (1,022 usable) is the smallest
fit; public needs 250, so /24 (254); data needs 100, so /25 (126);
management needs 20, so /27 (30). The total consumed is 10.20.0.0/22,
10.20.4.0/24 and the two halves of 10.20.5.0/24, leaving space for
growth.

\texttt{network/nacl.json} is the public-subnet NACL. Inbound rule 100
denies 198.51.100.0/24 for all protocols and ports before the web
allows; rule 110 allows TCP 80, rule 120 allows TCP 443, rule 130 allows
TCP 22 only from 10.20.5.128/27, and rule 140 allows return traffic from
10.20.0.0/22 on ephemeral ports 1024-65535. Outbound rule 100 allows TCP
8000 to 10.20.0.0/22, and rule 110 allows replies to clients on ports
1024-65535; there is an implicit deny for everything else. The official
simulator reports \textbf{6 / 6 = 100.0\%}.

A NACL is stateless: it does not remember that an inbound connection
created a flow, so the response packets are evaluated as new packets and
need their own rules in both directions. Rule order matters because NACL
rules are evaluated from the lowest number and the first match wins; the
deny for 198.51.100.0/24 must be numbered below the TCP 80/443 allows,
otherwise those allow rules would match first and the abusive range
would be allowed.

\hypertarget{task-6-reproducibility-testing-and-cicd}{%
\subsection{7. Task 6: reproducibility, testing and
CI/CD}\label{task-6-reproducibility-testing-and-cicd}}

\texttt{make\ up} creates the cluster, secrets, SeaweedFS, buckets,
images, uploads the model and applies the
Deployment/Service/Ingress/HPA; \texttt{make\ down} deletes the cluster.
From a fresh clone the whole system is two commands:
\texttt{make\ install-ml} and \texttt{make\ up}, or just
\texttt{make\ up} because the target bootstraps the virtual environment
if it is missing. \texttt{make\ up} does not commit credentials:
\texttt{scripts/create\_secrets.sh} generates them, writes them to the
git-ignored \texttt{.env.local} with mode 0600, and creates the
Kubernetes Secrets.

The test suite has 49 tests covering the cleaner's rejection precedence
and duplicate rule, accounting invariant, sharding and resume, subnet
arithmetic, the official NACL simulator, service API and exact error
body, readiness before load, output validation, storage backends, and
cost arithmetic. \texttt{ruff\ check\ .} passes, \texttt{pytest\ -q}
reports \texttt{49\ passed}, and \texttt{gitleaks\ detect\ -\/-no-git}
reports \texttt{no\ leaks\ found} (evidence in \texttt{results/}).

The CI workflow in \texttt{.github/workflows/ci.yml} runs in the
required order: \texttt{lint} (ruff) -\textgreater{} \texttt{unit-tests}
(pytest) -\textgreater{} \texttt{model-evaluation} (downloads
\texttt{model.tar.gz} and \texttt{model.sha256} from the \texttt{final}
GitHub Release, verifies the checksum, runs
\texttt{training/predict.py}, then \texttt{tools/score.py})
-\textgreater{} \texttt{quality-gate} (\texttt{-\/-min-validity\ 95}). A
separate \texttt{secret-scan} job runs gitleaks. A manual
\texttt{workflow\_dispatch} input \texttt{bad\_model=true} runs the
untrained Pythia-70M through the same gate and is expected to fail.

The CI run history for this repository, including the normal run and the
\texttt{bad\_model=true} quality-gate demonstration, is in its GitHub
Actions tab.

\texttt{AI\_USE.md} lists the AI tools and the checks applied; every
generated file was reviewed, linted, and tested, and the measured
results come from commands in \texttt{results/}, not from the assistant.

\hypertarget{task-7-performance-cost-and-review}{%
\subsection{8. Task 7: performance, cost and
review}\label{task-7-performance-cost-and-review}}

The local 50-request benchmark measured p50 \textbf{267.48 ms}, p95
\textbf{406.08 ms}, throughput \textbf{3.386 requests/s}, 50 total
requests, and 0\% errors. Ingress Locust results are 76 requests, 0
failures, p50 \textbf{6.0 s}, p95 \textbf{6.6 s}, and \textbf{1.297
requests/s}. The cold-start median is \textbf{11 s}. All results are in
\texttt{results/}; the short CPU-bound run did not demonstrate balanced
replica use.

\hypertarget{cost}{%
\subsubsection{8.1 Cost}\label{cost}}

A replica is billed for its CPU limit and memory limit for every hour it
runs, so with two replicas at 1 vCPU and 1 GiB each:

\begin{verbatim}
per replica per hour = 1 * 3.50 + 1 * 0.45 = 3.95 INR
all replicas per hour = 2 * 3.95 = 7.90 INR
requests per hour = measured requests/second * 3600
cost per 1,000 = (7.90 / requests per hour) * 1000
\end{verbatim}

Using the short Ingress run (1.297 requests/s, 76 responses, 0 HTTP
failures), the derived cost is approximately \textbf{1.69 INR per 1,000
successful HTTP responses}, assuming two replicas billed at their limits
for the entire hour. This is an estimate from a short run, not a
sustained capacity or semantic correctness measurement; Locust counted
HTTP success, not frame accuracy.

\texttt{scripts/cost.py} performs exactly this calculation from
\texttt{results/loadtest\_stats.csv} and \texttt{data/rate\_card.json},
and writes \texttt{results/cost.json}.

\hypertarget{slo}{%
\subsubsection{8.2 SLO}\label{slo}}

SLO: \textbf{99.5\% of \texttt{/parse} requests succeed (HTTP 200) over
one load-test window}. For the cluster window, N=76 and the error budget
is \texttt{floor(0.005\ *\ 76)\ =\ 0} failures. Actual failures were 0,
so that measured window met the SLO. This does not certify sustained
production availability.

\hypertarget{review}{%
\subsubsection{8.3 Review}\label{review}}

\begin{longtable}[]{@{}
  >{\raggedright\arraybackslash}p{(\columnwidth - 6\tabcolsep) * \real{0.2500}}
  >{\raggedright\arraybackslash}p{(\columnwidth - 6\tabcolsep) * \real{0.2500}}
  >{\raggedright\arraybackslash}p{(\columnwidth - 6\tabcolsep) * \real{0.2500}}
  >{\raggedright\arraybackslash}p{(\columnwidth - 6\tabcolsep) * \real{0.2500}}@{}}
\toprule\noalign{}
\begin{minipage}[b]{\linewidth}\raggedright
Task
\end{minipage} & \begin{minipage}[b]{\linewidth}\raggedright
Measured result
\end{minipage} & \begin{minipage}[b]{\linewidth}\raggedright
Limitation
\end{minipage} & \begin{minipage}[b]{\linewidth}\raggedright
Improvement
\end{minipage} \\
\midrule\noalign{}
\endhead
\bottomrule\noalign{}
\endlastfoot
1. Storage and labelling & 3,000 labels; 2,782 kept; accounting/coverage
pass & SeaweedFS exhausted volumes during v2 upload & Restore capacity
and verify v2 upload \\
2. Student training & 1,500 examples/version; v2 JSON validity 100\%,
full-frame 38\% vs teacher 40\% & Only 100 validation examples & Expand
held-out validation set \\
3. Service and container & 1.63 GB image; Kubernetes inference and local
50/50 pass & 1 CPU/1 GiB container benchmark not measured & Repeat
constrained benchmark with cgroup metrics \\
4. Kubernetes & Ingress 76 requests/0 failures; recovery 12 s;
cold-start median 11 s & Short run used one replica; v1-only update &
Longer test and true v1-to-v2 rollout \\
5. Networking & NACL simulator 6/6 correct & Simulated packets, not a
cloud deployment & Validate rules in the target cloud \\
6. Testing and CI & 49 tests pass; CI and expected quality-gate failure
recorded & No dedicated multi-node CI runtime & Add a disposable cluster
integration job \\
7. Performance and cost & p95 6.6 s; 1.297 RPS; estimated 1.69 INR/1,000
HTTP 200 & Short-run estimate, not semantic success cost & Measure
sustained throughput and correct frames \\
\end{longtable}

\hypertarget{limitations-and-improvements}{%
\subsection{9. Limitations and
improvements}\label{limitations-and-improvements}}

The main operational limitation is SeaweedFS volume exhaustion: the
service reported \texttt{no\ free\ volumes\ left} while uploading the v2
archive. The live deployment and Kubernetes behavior were therefore
validated with v1, while a real model artifact switch remains
unverified. The short Locust run also sent all recorded successful
requests to one replica, so it does not demonstrate balanced request
distribution. The image was built and run, but the separate 1 CPU / 1
GiB Docker benchmark was not measured.

Beyond the environment, a 70M-parameter student will not match a 3.8B
teacher; despite full 1,500-record training, the 100-example validation
set is small and the student's full-frame score remains below the
teacher. The HPA reacts to CPU rather than to queue depth, the service
serialises generation under a lock, and the k3d cluster is two nodes on
one machine. Better choices would be a queue-depth or latency-based
autoscaler, a batched/streaming inference server, and a multi-machine
cluster to test real node failure. On the data side, more teacher
examples, stricter slot grounding, and a held-out split would reduce
overfitting. On the security side, the S3 credentials are generated and
Secret-backed, the image is non-root and read-only, and the NACL is
default deny; a network policy restricting the app namespace to
SeaweedFS and the Ingress would add defence in depth.

\end{document}
